\section{The market economy in discrete time}\label{sec:qmkt:model}

The simulations are run on the market version rather than on the planner's, because the questions the
paper asks --- how long imperfect regulation delays the transition, and what the delay costs --- are
questions about policy instruments rather than about shadow prices. This part restates
Part~\ref{part:mkt} in discrete time, records the policy rates that implement the first best under the
timing conventions of Section~\ref{sec:qsp:setup:timing}, and sets out the scenarios and the welfare
measure.

\subsection{Agents in discrete time}\label{sec:qmkt:model:agents}

The market structure is that of Section~\ref{sec:mkt:setup:agents} unchanged: a household owning capital
and firms; competitive final-goods, mining and recycling sectors; and a waste-management sector under a
collection obligation with its subsidy set by tender. Prices are $P^R_t$, $P^T_t$ and $P^K_t$, with the
final good as numéraire, and the instruments are the seven of
Section~\ref{sec:mkt:setup:instruments}: $\tau^U_t$, $\tau^T_t$, $s^R_t$, $\tau^R_t$, $s^I_t$,
$\tau^{\mathcal D}_t$ and $\tau^{\circ}_t$, with $s^W_t$ pinned by the tender and $\mathcal T_t$
balancing the budget.

\smalltitle{Firms.} The static problems are unchanged, so their conditions carry over date by date:
\begin{align}
F_{K,t}\left(1-\tau^{\mathcal D}_t\Omega^Y_t\right) &= P^K_t,
&
F_{R,t}\left(1-\tau^{\mathcal D}_t\Omega^Y_t\right) &= P^R_t+\tau^R_t,
\label{eq:qmkt:final-foc}
\\
\alpha_t\left(P^R_t+s^R_t\right) &= P^T_t,
&
\frac{a^m_t}{\left(1+x_t\right)^2}\left(P^R_t+s^R_t\right) &= P^K_t
\quad\text{or}\quad K^R_t=0,
\label{eq:qmkt:recycler-foc}
\end{align}
the recycling corner binding when $a^m_t\left(P^R_t+s^R_t\right)\le P^K_t$, and the waste firm choosing
$\varpi_t$ from the three-branch condition
\begin{align}
P^T_t+\tau^U_t-\tau^T_t &= \frac{\bar c^T}{\left(1+g^T\right)^t}\varpi_t^{\xi^T}
\left(1+\tau^{\mathcal D}_t\Omega^W_t\right)
\label{eq:qmkt:waste-foc}
\end{align}
on the interior branch, with $\varpi_t=0$ and $\varpi_t=1$ at the two corners of
Section~\ref{sec:mkt:setup:firms}, and with
\begin{align}
s^W_t &= c^W_t\left(1+\tau^{\mathcal D}_t\Omega^W_t\right)+\tau^U_t\left(1-\varpi_t\right)
+\tau^T_t\varpi_t-P^T_t\varpi_t
\label{eq:qmkt:waste-zero-profit}
\end{align}
from the tender. The mining firm now discounts at the market rate, so~\eqref{eq:mkt:mining-costates}
becomes the pair of recursions
\begin{align}
P^S_t &= \left|C^N_{S,t}\right|\left(1+\tau^{\mathcal D}_t\Omega^N_t\right)+\hat P^S_t,
&
P^X_t &= -C^D_{X,t}\left(1+\tau^{\mathcal D}_t\Omega^D_t\right)+\hat P^X_t ,
\label{eq:qmkt:mining-costates}
\end{align}
with $\hat P_t\equiv P_{t+1}/\left(1+r_{t+1}\right)$ as in~\eqref{eq:qsp:discount}, and its two flow
conditions
\begin{align}
P^R_t &= C^N_{N,t}\left(1+\tau^{\mathcal D}_t\Omega^N_t\right)+\hat P^S_t+\tau^{\circ}_t\Omega^{N,S}_t,
&
\hat P^S_t+\hat P^X_t &= C^D_{D,t}\left(1+\tau^{\mathcal D}_t\Omega^D_t\right)+\tau^{\circ}_t\Omega^{D,S}_t,
\label{eq:qmkt:mining-foc}
\end{align}
each with the corner branch that shuts the activity down.

\smalltitle{Household.} The budget and the two states are
\begin{align}
\left(1+\tau^{\mathcal D}_t\Omega^C_t\right)C_t+\left(1-s^I_t\right)G_t+\tau^{\circ}_t\delta M^K_t
&= P^K_tK_t+\Pi_t+\mathcal T_t,
\label{eq:qmkt:hh-budget}
\end{align}
with $K_{t+1}=\left(1-\delta\right)K_t+G_t$ and $M^K_{t+1}=\left(1-\delta\right)M^K_t+\Phi^I_tG_t$, and
the conditions are
\begin{align}
\frac{C_t^{-\eta}}{\lambda_t} &= 1+\tau^{\mathcal D}_t\Omega^C_t,
&
\hat q_t &= 1-s^I_t+\Phi^I_t\hat p^M_t,
&
p^M_t &= \delta\tau^{\circ}_t+\left(1-\delta\right)\hat p^M_t,
\label{eq:qmkt:hh-foc}
\\
q_t &= P^K_t+\left(1-\delta\right)\hat q_t,
\label{eq:qmkt:hh-arbitrage}
\end{align}
together with the Euler equation implied by~\eqref{eq:qsp:discount}. The household takes $\Omega^C_t$
and $\Phi^I_t$ as given, which is where the intensity externality of
Section~\ref{sec:mkt:setup:externalities} sits, and it takes $P_t$ as given, which is where the
pollution externality sits.

\subsection{The implementing policy}\label{sec:qmkt:model:policy}

\begin{proposition}[Discrete-time implementation]\label{prop:qmkt:implementation}
Let $\left\{p^P_t,p^W_t,\zeta_t,\Phi^I_t,r_t\right\}$ solve Problem~\ref{prob:qsp}, and set
\begin{align}
\tau^{\Xi}_t &= \hat p^P_t\;=\;\frac{p^P_{t+1}}{1+r_{t+1}},
&
\tau^{\mathcal D}_t &= \zeta_t,
&
\tau^{\circ}_t &= p^W_t,
\label{eq:qmkt:implement-primitives}
\end{align}
so that $\tau^U_t=\hat p^P_t$, $\tau^T_t=s^R_t=d^W\hat p^P_t$, $\tau^R_t=p^W_t-\zeta_t$ and
$s^I_t=\Phi^I_t\left(p^W_t-\zeta_t\right)$. Then the competitive equilibrium implements the first best;
the equilibrium prices are $P^K_t=F_{K,t}\left(1-\zeta_t\Omega^Y_t\right)$, $P^R_t=\Psi_t-p^W_t$,
$P^T_t=\alpha_tV_t$ and $s^W_t=p^W_t$, with $P^S_t=p^S_t$ and $P^X_t=p^X_t$; and net government revenue
is $\hat p^P_t\Xi_t$ at every date.
\end{proposition}

\begin{proof}
Every step of the proof of Proposition~\ref{prop:mkt:implementation} goes through with $p^P$ replaced
by $\hat p^P_t$ and with each stock price replaced by its present value, which is precisely the
correspondence established in Section~\ref{sec:qsp:model:lagrangian}. The two price identities
$\Psi_t=P^R_t+\tau^R_t+\tau^{\mathcal D}_t=P^R_t+p^W_t$ and $V_t=P^R_t+s^R_t$ are unchanged, the
recycling and treatment branches match term for term, the zero-profit
condition~\eqref{eq:qmkt:waste-zero-profit} returns $s^W_t=p^W_t$ by the same cancellation, and the
household's $\hat q_t$ in~\eqref{eq:qmkt:hh-foc} reproduces~\eqref{eq:qsp:foc:G} once $s^I_t$ is
substituted. The budget identity is the computation of Corollary~\ref{cor:mkt:budget} with $\hat p^P_t$
in place of $p^P$.
\end{proof}

\smalltitle{The one thing discretisation changes.} The Pigouvian rate is $\hat p^P_t$, not $p^P_t$: a
tonne emitted in period $t$ enters the pollution stock only in $t+1$, so it is charged the
\emph{discounted next-period} shadow cost. In continuous time the distinction vanishes; in a model
calibrated at annual frequency with a long horizon it does not, and getting it wrong biases the
emission tax upward by a factor $1+r$. The materials instruments are unaffected, because waste is
collected, treated and recycled within the period in which it is generated.

\subsection{Scenarios}\label{sec:qmkt:scenarios}

Because the instrument set separates cleanly into three blocks, the natural experiments are the
$2^3-1$ ways of switching blocks off, plus the technology experiments. The following are the ones worth
running first.

\begin{enumerate}[label=\emph{(\arabic*)},leftmargin=2.4em,itemsep=0.2em]
\item \emph{Benchmark.} Full optimal policy~\eqref{eq:qmkt:implement-primitives} and neutral technical
progress~\eqref{eq:qsp:neutral-progress}. This reproduces the planner's path and dates $t_R$.

\item \emph{Laissez-faire.} Every instrument zero, as in Section~\ref{sec:mkt:implementation:unregulated}.
The materials side of the model goes silent, $p^M_t=0$ and $q_t=1$, and the only remaining incentive to
recycle is the market price of material. This is the upper bound on the delay.

\item \emph{Emission tax only.} $\tau^\Xi_t=\hat p^P_t$ with $\tau^{\mathcal D}_t=\tau^{\circ}_t=0$,
hence $\tau^R_t=s^I_t=0$. This is the policy the earlier literature would prescribe --- and the
comparison with~(1) is the paper's cleanest statement of what the materials accounting adds. It is the
scenario most directly comparable to \textcite{sorensen2018}, where a pollution tax alone is sufficient
because extraction is the only source of pollution.

\item \emph{Materials block only.} $\tau^{\mathcal D}_t=\zeta_t$ and $\tau^{\circ}_t=p^W_t$ with
$\tau^\Xi_t=0$. The mirror image of~(3), and the one that isolates how much of the transition is driven
by the materials ledger rather than by environmental damage.

\item \emph{No recycling instruments.} $\tau^U_t$, $\tau^T_t$ and $s^R_t$ set to zero with the rest at
their optimal levels. This is the scenario in the paper draft's list, and it isolates the contribution
of waste treatment and recycling policy specifically.

\item \emph{Shutting the composition channel.} $\phi^I=0$, which by~\eqref{eq:qsp:foc:Omega} sets
$\zeta_t\equiv0$ and makes $\tau^{\mathcal D}_t=0$ and $\tau^R_t=\tau^{\circ}_t=p^W_t$. This is not a
policy experiment but a decomposition: it measures how much of the model's behaviour comes from the
composition of final-good uses as against the ledger alone, and it is the natural robustness check on
the part of the framework that is hardest to calibrate. The appendix variants of
Sections~\ref{sec:app:exogenousPhiI} and~\ref{sec:app:constantPhiI} belong in the same family.

\item \emph{Sectoral technical progress.} Departures from~\eqref{eq:qsp:neutral-progress} one rate at a
time. Remark~\ref{rem:qsp:tail} predicts that $g^R$ dominates, since it is the only rate that relaxes
the cap on the circularity multiplier; $g^N$ should work in the opposite direction by keeping virgin
extraction cheap and postponing $t_R$.

\item \emph{The substitution limit.} $\bar R$ raised from zero, per Remark~\ref{rem:qsp:Rbar}. The
interesting question is whether a binding floor advances the transition, by making material scarcity
bite earlier, or renders it infeasible, by requiring a circularity multiplier the frontier $a^m_t$
cannot deliver.
\end{enumerate}

For each scenario the objects to report are the transition date $t_R$, the paths of $\varpi_t$ and
$K^R_t/K_t$ --- whose ordering is the substantive conjecture of
Section~\ref{sec:sp:opt:results} --- the date at which $P^T_t$ turns positive, which by
Corollary~\ref{cor:mkt:phases} coincides with $t_R$ under optimal policy and need not otherwise, the
material throughput $R_t$ and the circularity multiplier $\left(1-a_t\varpi_t\right)^{-1}$, the signs
and crossings of $p^W_t$, $p^M_t$ and $\zeta_t$, and the welfare loss defined next.

\subsection{Welfare}\label{sec:qmkt:welfare}

Split lifetime utility into its two components,
\begin{align}
U_0 &= U^C_0-U^P_0,
&
U^C_0 &= \sum_{t\ge0}\beta^t\frac{C_t^{1-\eta}}{1-\eta},
&
U^P_0 &= \sum_{t\ge0}\beta^t\frac{\psi_t}{1+\nu}P_t^{1+\nu},
\label{eq:qmkt:welfare-split}
\end{align}
and write $U^{C*}_0$ and $U^{P*}_0$ for their values along the first-best path. Define the welfare cost
of a policy as the permanent proportional increase $\vartheta$ in consumption, applied to every period
of the scenario path, that would leave the household as well off as under the optimal policy:
\begin{align}
\left(1+\vartheta\right)^{1-\eta}U^C_0-U^P_0 &= U^{C*}_0-U^{P*}_0
\qquad\Longleftrightarrow\qquad
\vartheta = \left[\frac{U^{C*}_0-U^{P*}_0+U^P_0}{U^C_0}\right]^{\frac{1}{1-\eta}}-1 .
\label{eq:qmkt:cev}
\end{align}
The closed form follows because the felicity function is homogeneous of degree $1-\eta$ in consumption
while the pollution term is separable, so scaling consumption scales $U^C_0$ and leaves $U^P_0$ alone.
The numerator makes the decomposition transparent: the compensation has to cover the extra disutility
from pollution, $U^P_0-U^{P*}_0$, together with any shortfall in the consumption stream itself.

This is a consumption-equivalent measure rather than the wealth-equivalent compensating variation used
in the paper draft. The two answer the same question and we expect them to be close; the
consumption-equivalent version is preferred here because it requires nothing but the two simulated
paths, whereas the wealth-based version requires the household's intertemporal budget constraint,
present-value profit and a pair of intertemporal price indices to be constructed consistently under
each policy --- more machinery, and more places for a policy-dependent price index to introduce an
artefact. \emph{Reporting both on the benchmark comparison is a sensible check, and is worth doing once
the wealth-based version is available.}

\subsection{Solving the market model}\label{sec:qmkt:solution}

Under the optimal policy, Proposition~\ref{prop:qmkt:implementation} makes the market model redundant:
solve Problem~\ref{prob:qsp} and read the instruments off the shadow prices. It is the \emph{other}
scenarios that require a separate algorithm, and they are harder than the planner's problem in one
specific respect. With instruments off their optimal values there is no planning problem whose solution
the equilibrium is, so the trajectory has to be found by solving the equilibrium conditions directly
rather than by optimising --- method~(a) of Section~\ref{sec:qsp:model:solution} is unavailable, and
the natural formulation is a large nonlinear system in the whole trajectory, or shooting on the four
private costates $\left(q_t,p^M_t,P^S_t,P^X_t\right)$.

Two practical points follow. First, the within-period block of Section~\ref{sec:qsp:model:block} is
unchanged by policy --- it is accounting, not optimisation --- so it can be built once and reused across
all scenarios. Second, the phase switches are now \emph{policy-dependent}: the recycling corner binds
when $a^m_t\left(P^R_t+s^R_t\right)\le P^K_t$, which under scenario~(5) has no $s^R_t$ in it and under
scenario~(2) has neither the subsidy nor the deposit that raises $P^R_t$. A robust corner-detection
routine is therefore not an optional refinement but the part of the implementation that the paper's
headline number --- how long imperfect regulation delays the transition --- depends on directly.
